\documentclass[11pt, letterpaper, includehead]{article}

%%%%%%%%%%%%%%%%%%%%% Pre-document %%%%%%%%%%%%%%%%%%%%%
\usepackage{fancyhdr}  % Allow for headers
\usepackage{graphicx}  % Allow for figures 
\usepackage{float}     % Allow for figure inserted in specified location
\usepackage{amsmath}   % Allow for aligned math
\usepackage{array}     % Allow for cell width manipulation
\usepackage{nicematrix}
\usepackage{amssymb} % Uhhhh what was this????
\usepackage{multicol}

\setlength{\parindent}{0pt} % Remove auto paragraph indents

% Get rid of those big ass margins
\usepackage[margin=1in]{geometry}

% Table cell formatting
\setlength{\arrayrulewidth}{0.25mm}
\setlength{\tabcolsep}{11pt}
\renewcommand{\arraystretch}{1.2}

\begin{document}

%%%%%%%%%%%%%%%%%%%%% Title Page %%%%%%%%%%%%%%%%%%%%%
\begin{titlepage}
  \begin{center}
    \Huge{\textbf{Lab 4}}\\
    \Huge{Projectile Motion}
    \vfill
    \begin{figure}[H] % H makes the figure insert at the position in the document
      \centering 
      \includegraphics[width=6cm]{../logo.png}
    \end{figure}
    \large{\textbf{your name here}}\\
    \large{Julian Barossi, Liam Gilligan, Stephanie L'Heureux}\\
    \vspace{0.5cm}
    \normalsize
    \today
  \end{center}
\end{titlepage}

%%%%%%%%%%%%%%%%%%%%% TABLE OF CONTENTS %%%%%%%%%%%%%%%%%%%%%
\tableofcontents
\pagebreak % Move to next page

% Add a nice fancy header
\pagestyle{fancy}
\fancyhead{}
\fancyhead[C]{\textbf{Lab 4:} Projectile Motion}

\section{Measuring the launch Speed (Horizontal Projections)} % 1

\subsection{Measurements} % 1.1

\subsubsection{Verticle distance} % 1.1.1
The distance $\Delta y$ for the ball to fall from the lab table was measured in two parts. 
First, we measured the height from the bottom of the ball to the bottom of the launcher on the lab table.
This distance was found to be $24.20cm$. Then the distance from the edge of the lab table to the 
floor was $93.70cm$. Combining these we get $\Delta y$.
$$\Delta y = 24.20cm + 93.70cm = 117.90cm \Rightarrow 1.18m$$ 
$$\boxed{\Delta y = 1.18m}$$ 

\subsubsection{Horizontal distance} % 1.1.2
We fired the ball from the lab table at a launch angle of $\theta = 0^{\circ}$ and recorded $\Delta x$.
\begin{center} 
  \begin{tabular}{|  m{2.2cm} | m{3.2cm} | m{3.2cm} | m{3.2cm} | } 
    \hline
    \textbf{Attempt} & \boldmath{$\Delta{x_{paper}} \, (cm)$} & \boldmath{$\Delta{x_{total}}\, (cm)$} & \boldmath{$\Delta{x} (m)$} \\ 
       \hline
     1 & 20.25 & 349.25 & 3.49 \\
       \hline
  \end{tabular} 
\end{center}

\subsection{Calculations} % 1.2
\begin{multicols}{3}
  \centering{\textbf{Time in air (t)}}
  $$\Delta y = V_{0_y}t + \frac{1}{2}a_yt^2$$
  $$\Delta y = \frac{1}{2}a_yt^2$$
  $$t = \sqrt{\frac{2\Delta y}{a_y}}$$
  $$t = \sqrt{\frac{2(-1.18m)}{-9.8 m/s^2}}$$
  $$\boxed{t = -0.491s}$$
  
  \columnbreak
  
  \centering{\boldmath{\textbf{Horizontal velocity} ($V_x$)}}
  $$\Delta x = V_{0_x}t + \frac{1}{2}a_xt^2$$
  $$\Delta x = V_{0_x}t$$
  $$V_{0_x} = \frac{\Delta x}{t}$$
  $$V_{0_x} = \frac{3.4925m}{0.491s}$$
  $$V_{0_x} = 7.113 \, m/s \approx 7.11\,m/s$$
  $$\boxed{V_{0_x} = 7.11\,m/s}$$

  \columnbreak

  \centering{\boldmath{\textbf{Launch speed} ($V$)}}
  $$V = \sqrt{V_x^2 + V_y^2}$$
  $$V = \sqrt{(7.113 \, m/s)^2 + (0 \, m/s)^2}$$
  $$V = 7.113 \, m/s \approx 7.11 \, m/s$$
  $$\boxed{V = 7.11 \, m/s}$$
\end{multicols}

\subsection{Data} % 1.3
\subsubsection{Average launch speed, standard deviation, and a standard error}
\begin{center} 
  \begin{tabular}{|  m{2.2cm} | m{2.2cm} | m{2.2cm} | m{2.2cm} | m{2.2cm} | } 
    \hline
    \textbf{Attempt} & \boldmath{$\Delta{x_{paper}} \, (cm)$} & \boldmath{$\Delta{x_{total}}\, (cm)$} & \boldmath{$\Delta{x} (m)$} & \boldmath{$ V_{0_x} \, (m/s)$} \\ 
       \hline
     1 & 20.25 & 349.25 & 3.49 & 7.11 \\ 
       \hline
     2 & 25.60 & 354.60 & 3.55 & 7.23 \\ 
       \hline
     3 & 22.50 & 351.50 & 3.52 & 7.17 \\ 
       \hline
     4 & 23.90 & 352.90 & 3.53 & 7.19 \\ 
       \hline
     5 & 20.00 & 349.00 & 3.49 & 7.11 \\ 
       \hline
     6 & 22.10 & 351.10 & 3.51 & 7.16 \\ 
       \hline
     7 & 19.40 & 348.40 & 3.48 & 7.10 \\ 
       \hline
     8 & 22.40 & 351.40 & 3.51 & 7.16 \\ 
       \hline
     9 & 21.90 & 350.90 & 3.51 & 7.15 \\ 
       \hline
     10 & 22.00 & 351.00 & 3.51 & 7.16 \\ 
       \hline
  \end{tabular} 
\end{center}

\begin{center} 
  \begin{tabular}{| m{3cm} |  m{6cm} |m{3.6cm} |} 
    \hline
     \textbf{Quantity} & \boldmath{\textbf{Value} ($m/s$)}  & \boldmath{\textbf{Rounded} ($m/s$)} \\ 
       \hline
        $ \bar{V}_{0_x}$ & 7.15573816177454 & 7.16 \\ 
       \hline
        $\sigma$ & 0.0377269809207384 & 0.038 \\ 
       \hline
        SE & 0.011930318895125 & 0.012 \\
       \hline
  \end{tabular} 
\end{center}

\subsubsection{Percentage of values which fall in one standard deviation}
$$\frac{\# measurements \in [\bar{V_{0_x}} - \sigma, \bar{V_{0_x}} + \sigma]}{N} \\$$
$$\frac{\# measurements \in [7.156\,m/s - 0.0377\,m/s, 7.156\,m/s + 0.0377\,m/s]}{10} \\$$
$$\frac{\# measurements \in [7.1183\,m/s, 7.1937\,m/s]}{10} \\$$
$$\boxed{70\%}$$

\subsubsection{Does the data match the statistical prediction} 
By statistical theory, the dispersal of data points and the standard deviation relate such that $68\%$ of 
data points fall within $\geq (\bar{V_{0_x}} - \sigma)$ and $\leq (\bar{V_{0_x}} + \sigma)$. Our data set nearly matches theory 
with $70\%$ of values falling within one standard deviation of the mean.

\subsection{Best value for initial speed} % 1.4
$$\bar{V} \pm 2SE$$
$$7.16\, m/s \pm 2( 0.012 \, m/s)$$
$$\boxed{7.16 \, m/s \pm 0.024 \, m/s}$$

\section {Measuring the launch speed (Vertical Projections)} % 2
\subsection{Max y height reached} % 2.1
\begin{center} 
  \begin{tabular}{|  m{3cm} | m{6cm} | m{3.6cm} |  } 
    \hline
    \textbf{Attempt} & \boldmath{$\Delta{y}\, (cm)$} & \boldmath{$\Delta{y} (m)$} \\ 
       \hline
        1 & 214.10 & 2.14 \\
       \hline
  \end{tabular} 
\end{center}

\subsection{Launch speed} % 2.2
By plugging in the displacement from the end of the spring
in the launcher to the maximum height into $\Delta y$, we
can calculate the launch speed for that measurement.

$$V_y^2 = V_{0_y}^2 + 2a_y\Delta y$$
$$0\,m/s = V_{0_y}^2 + 2a_y\Delta y$$
$$-V_{0_y}^2 = 2a_y\Delta y$$
$$V_{0_y}^2 = -2a_y\Delta y$$
$$V_{0_y} = \sqrt{-2a_y\Delta y}$$
$$V_{0_y} = \sqrt{-2(-9.8\, m/s^2)(2.141\, m )}$$
$$\boxed{V_{0_y} = 6.48 \, m/s}$$

\subsection{Data} % 2.3
\subsubsection{Average launch speed, standard deviation, and a standard error} % 2.3
\begin{center} 
  \begin{tabular}{|  m{2.2cm} | m{3.2cm} | m{3.2cm} | m{3.2cm} | }
    \hline 
    \textbf{Attempt} & \boldmath{$\Delta{y} \, (cm)$} & \boldmath{$\Delta{y}\, (m)$} & \boldmath{$V_{0_y} \, (m/s)$} \\ 
    \hline
     1 & 214.10 & 2.14 & 6.48 \\ 
       \hline
     2 & 214.80 & 2.15 & 6.49 \\ 
       \hline
     3 & 216.10 & 2.16 & 6.51 \\ 
       \hline
     4 & 216.70 & 2.17 & 6.52 \\ 
       \hline
     5 & 217.40 & 2.17 & 6.53 \\ 
       \hline
     6 & 216.90 & 2.17 & 6.52 \\ 
       \hline
     7 & 217.60 & 2.18 & 6.53 \\ 
       \hline
     8 & 217.50 & 2.18 & 6.53 \\ 
       \hline
     9 & 219.70 & 2.20 & 6.56 \\ 
       \hline
     10 & 221.10 & 2.21 & 6.58 \\ 
       \hline
  \end{tabular} 
\end{center}

\begin{center} 
  \begin{tabular}{| m{3cm} |  m{6cm} |m{3.6cm} |} 
    \hline
     \textbf{Quantity} & \boldmath{\textbf{Value} ($m/s$)}  & \boldmath{\textbf{Rounded} ($m/s$)} \\ 
       \hline
        $ \bar{V}_{0_y}$ & 6.52444397275332 & 6.52 \\ 
       \hline
        $\sigma$ & 0.0311063838193513 & 0.031 \\ 
       \hline
        SE & 0.00983670226405577 & 0.0098 \\
       \hline
  \end{tabular} 
\end{center}

\subsubsection{Percentage of values which fall in one standard deviation}
$$\frac{\# measurements \in [\bar{V_{0_y}} - \sigma, \bar{V_{0_y}} + \sigma]}{N} \\$$
$$\frac{\# measurements \in [6.524\,m/s - 0.0311\,m/s, 6.524\,m/s + 0.0311\,m/s]}{10} \\$$
$$\frac{\# measurements \in [6.4929\,m/s, 6.5551\,m/s]}{10} \\$$
% 68100\frac{68}{100}
$$\boxed{60\%}$$

\subsubsection{Does the data match the statistical prediction} 
As previously discussed, $68\%$ of 
data points fall within $\geq (\bar{V_{0_y}} - \sigma)$ and $\leq (\bar{V_{0_y}} + \sigma)$. 
For this data set, $60\%$ fell within this range. This inconstancy is likely due to the 
low quantity of samples in the set. 

\subsection{Best value for initial speed} % 1.4
$$\bar{V} \pm 2SE$$
$$6.52\, m/s \pm 2( 0.0098 \, m/s)$$
$$\boxed{6.52\, m/s \pm 0.020 \, m/s}$$

\section{Comparing your Results} % 3
\textbf{Which is most precise? Does this mean that it is the most accurate?}\\ 

\textbf{Are the measured vertical and horizontal launch speeds the same or is there a
significant difference between them? Should they be identical in theory? Why or
why not?}\\

\textbf{What value for initial speed should you use for making predictions in part 4, and
why?}
$$V = (0.8 * V_{0_x}) + (0.2 * V_{0_y})$$
$$V_0 = (0.8 * 7.16\,m/s) + (0.2 * 6.52\,m/s)$$
$$V_0 = 7.032\,m/s \approx 7.03\,m/s$$

\section{Challenge Shot} % 4 
\subsection{Calculations} % 4.1
\subsubsection{Optimal theta} % 4.1.1
Rather than using the suggested launch angle of $\theta = 45^{\circ}$, we choose
to find an "optimal theta", an angle for which the ball would land in the 
target at its maximum height, where $V_{0_y} = 0\,m/s$. This would prevent us from 
overshooting the target. For $\Delta y$, we averaged the 
distances of top pole and and bottom pole. $\Delta y = 0.48m$.

$$V_y^2 = V_{0_y}^2 + 2a_y\Delta y$$
$$0\,m/s = V_{0_y}^2 + 2a_y\Delta y$$
$$-V_{0_y}^2 = 2a_y\Delta y$$
$$V_{0_y}^2 = -2a_y\Delta y$$
$$V_{0_y} = \sqrt{-2a_y\Delta y}$$
$$V_0\cdot\sin\theta = \sqrt{-2a_y\Delta y}$$
$$\sin\theta = \frac{\sqrt{-2a_y\Delta y}}{V_0}$$
$$\theta = \arcsin\left(\frac{\sqrt{-2a_y\Delta y}}{V_0}\right)$$
$$\theta = \arcsin\left(\frac{\sqrt{-2(-9.80m/s^2)(0.4845m)}}{7.03m/s}\right)$$
$$\theta = \arcsin(0.4377)$$
$$\boxed{\theta = 26^{\circ}}$$

\subsubsection{Time to maximum height} % 4.1.2
$$V_y = V_{0_y} + a_yt$$
$$0\,m/s = V_{0_y} + a_yt$$
$$-V_{0_y} = a_yt$$
$$t = \frac{-V_{0_y}}{a_y}$$
$$t = \frac{-7.03m/s\cdot\sin26^{\circ}}{-9.80m/s^2}$$
$$t = 0.314441s \approx 0.31s$$
$$\boxed{t = 0.31s}$$

\subsubsection{Horizontal distance} % 4.1.3
$$\Delta x = V_{0_x}t + \frac{1}{2}a_xt$$
$$\Delta x = V_{0_x}t$$
$$\Delta x = 7.03m/s\cdot \cos26^{\circ}\cdot 0.31s$$
$$\Delta x = 1.98666m \approx 1.99m$$
$$\boxed{\Delta x = 1.99m}$$
Based on the calculations, we positioned the launcher at $26^{\circ}$, $1.99m$ from the center
of the target.

\subsection{Analysis}  % 4.2
Our calculations proved to be correct. Our first shot landed within the target.
We believe choosing an optimal angle theta, rather than the suggested value of 
$45^{\circ}$, improved our chances of landing in the target.

\end{document}
